% ============================== Part B ==============================
\problem{B1}{Bias, scatter and pulls}

\begin{given}
An estimator $\hat\theta$ of a true value $\theta$ comes with a reported error bar $\hat\sigma$. The \emph{pull} is
$(\hat\theta - \theta)/\hat\sigma$. Suppose $\hat\theta \sim \mathcal N(\theta + b, s^2)$: bias $b$, true scatter $s$.
\emph{Derived here}; the pull is a standard diagnostic.
\end{given}

\need{nothing beyond Gaussians.}

\begin{enumerate}
  \item With $\hat\sigma$ fixed, show that the pull is $\mathcal N(b/\hat\sigma,\ s^2/\hat\sigma^2)$.
  \item What do ``pull mean $\neq 0$'' and ``pull spread $\neq 1$'' each tell you?
  \item Our 50-realization ensemble gives a pull spread of 0.84. By what factor do the reported error bars exceed
    the true scatter? If you used them in an inference, what fraction of the available information
    ($\propto 1/\sigma^2$) would you be claiming?
  \item The polar cells give 1437 pulls. What is the statistical uncertainty on their mean? Is a pull mean of
    $-0.02$ evidence of bias?
\end{enumerate}

\solution{
(a) $(\hat\theta - \theta)/\hat\sigma$ is a linear function of a Gaussian, with mean $b/\hat\sigma$ and standard
deviation $s/\hat\sigma$.
(b) A non-zero mean means bias, in units of the error bar. A spread above 1 means the error bars are too small
(overconfident); below 1, too large (conservative).
(c) $\hat\sigma/s = 1/0.84 = 1.19$, so the errors are 19\% too large. You would claim
$s^2/\hat\sigma^2 = 0.84^2 = 0.71$ of the information, as if you had thrown away 29\% of the galaxies. This loses
precision but introduces no bias.
(d) $0.84/\sqrt{1437} = 0.022$. A mean of $-0.02$ is $0.9\sigma$: consistent with zero.
}

% --------------------------------------------------------------------
\problem{B2}{Multiplicative and additive bias; shrinkage of a posterior mean}

\begin{given}
Shear biases are summarised by a linear response $\hat g = (1+m)\,g + c$, with $m$ the multiplicative and $c$ the
additive bias \source{\cBAresponse.} Model one galaxy with a Gaussian likelihood $d \sim \mathcal N(g, \sigma^2)$ and a
Gaussian prior $g \sim \mathcal N(0, \sigma_\pi^2)$. \emph{Derived here} (it reproduces the toy model of paper
Sec.~6.2).
\end{given}

\need{A1 (product of two Gaussians).}

\begin{enumerate}
  \item Treat the prior as a second ``measurement'' $0 \pm \sigma_\pi$. Show that the posterior mean is
    $\hat g = d\,\sigma_\pi^2/(\sigma_\pi^2 + \sigma^2)$.
  \item Average over the noise to find $\E[\hat g]$ and read off $m$ and $c$.
  \item Use $\sigma = 0.07$ and the standard deviation of a uniform prior on $[-0.1, 0.1]$,
    $\sigma_\pi = 0.2/\sqrt{12}$, to estimate $m$. Compare with A5(c).
\end{enumerate}

\solution{
(a) By A1, with weights $1/\sigma^2$ on $d$ and $1/\sigma_\pi^2$ on $0$:
$\hat g = (d/\sigma^2)/(1/\sigma^2 + 1/\sigma_\pi^2) = d\,\sigma_\pi^2/(\sigma_\pi^2+\sigma^2)$.
(b) $\E[d] = g$, so $\E[\hat g] = g\,\sigma_\pi^2/(\sigma_\pi^2 + \sigma^2)$: $m = -\sigma^2/(\sigma^2 + \sigma_\pi^2)$
and $c = 0$.
(c) $\sigma_\pi = 0.0577$, so $m = -0.0049/(0.0049 + 0.00333) = -0.60$. The truncated box of A5 gave $-0.57$; the
Gaussian toy captures the effect well. Shrinkage is strongest when the likelihood is wide compared with the prior.
}

% --------------------------------------------------------------------
\problem{B3}{Stack likelihoods, not posteriors}

\begin{given}
Each galaxy's NPE posterior is $q_i(g) \propto L_i(g)\,\pitrain(g)$ (A4). $N$ galaxies share one $g$ and have
independent data. The correct joint posterior is $p(g\mid d_1,\dots,d_N) \propto \pi(g)\prod_i L_i(g)$, and samples
drawn under an interim prior are combined by dividing that prior out.
\source{\cBAconst; \cInterim.}
\end{given}

\need{A4, B2.}

\begin{enumerate}
  \item Show that $\prod_i q_i(g) \propto \pitrain(g)^N\prod_i L_i(g)$. How many times is the prior counted?
  \item Use the Gaussian model of B2 (prior $\mathcal N(0,\sigma_\pi^2)$, each likelihood width $\sigma$). Show that the
    mean of $\prod_i q_i$ is shrunk by the same factor $\sigma_\pi^2/(\sigma_\pi^2+\sigma^2)$ for every $N$.
  \item What is the mean of the correct posterior $\pi\prod_i L_i$, and how does its $m$ behave as $N \to \infty$?
  \item Why does our stack use $L_i = q_i/\pitrain$ and a single flat prior on the cell shear?
\end{enumerate}

\solution{
(a) $\prod_i q_i \propto \prod_i L_i\pitrain = \pitrain^N\prod_i L_i$: the prior is counted $N$ times.
(b) $\pitrain^N$ has precision $N/\sigma_\pi^2$ and $\prod L_i$ has precision $N/\sigma^2$ about the mean $\bar d$, so
the mean is $\bar d\,(N/\sigma^2)/(N/\sigma^2 + N/\sigma_\pi^2) = \bar d\,\sigma_\pi^2/(\sigma_\pi^2+\sigma^2)$: the
$N$ cancels, and $m = -0.6$ no matter how many galaxies you multiply.
(c) $\bar d\,(N/\sigma^2)/(N/\sigma^2 + 1/\sigma_\pi^2)$, so $m = -\sigma^2/(\sigma^2 + N\sigma_\pi^2) \to 0$ as
$N \to \infty$.
(d) Dividing out $\pitrain$ leaves the likelihoods, which multiply correctly. The cell then has one prior (uniform
on $|\vG| < 0.1$), counted once, and the data dominate as $N$ grows.
}

% --------------------------------------------------------------------
\problem{B4}{Fitting one shear to galaxies whose shears differ: the reference shear}

\begin{given}
Fitting a model that cannot be exactly right (a misspecified model) converges, as the data grow, to the parameter
value that makes the model's distribution closest to the true one in Kullback--Leibler divergence: the
\emph{pseudo-true} value.
\source{\cMisspec.}
In a cell, galaxy $i$ has true shear $\vt_i$ (in the cell frame), a Gaussian likelihood centred on $\vd_i$ with
precision $F_i$, and $\E[\vd_i] = \vt_i$. The stack fits one common $\vG$ to all of them.
\end{given}

\need{A1, A3 (the matrix-weighted mean).}

\begin{enumerate}
  \item By A3, $\hat\vG = \Ftot^{-1}\sum_i F_i\vd_i$. Average over the noise, assuming $F_i$ doesn't depend on it.
    Show that $\E[\hat\vG] = \Ftot^{-1}\sum_i F_i\vt_i \equiv \gref$.
  \item Show that $\mathrm{Cov}(\hat\vG) = \Ftot^{-1}$, and that it doesn't depend on how much the $\vt_i$ differ.
    Why is this still the right error bar \emph{about} $\gref$?
  \item Scalar example: galaxy 1 has $t_1 = 0.02$, $\sigma_1 = 0.02$; galaxy 2 has $t_2 = 0$, $\sigma_2 = 0.06$.
    Compute $g_{\rm ref}$ and the plain mean. Which one does the stack estimate?
  \item Why would comparing the stack with the shear at the cell centre be the wrong test?
\end{enumerate}

\solution{
(a) $\E[\hat\vG] = \Ftot^{-1}\sum_i F_i\,\E[\vd_i] = \Ftot^{-1}\sum_i F_i\vt_i$. This is the pseudo-true value for
Gaussian likelihoods: the constant that best describes the varying $\vt_i$, with each galaxy weighted by its
precision.
(b) With independent noises of covariance $F_i^{-1}$:
$\mathrm{Cov}(\hat\vG) = \Ftot^{-1}\big(\sum_i F_iF_i^{-1}F_i\big)\Ftot^{-1} = \Ftot^{-1}$. The $\vt_i$ don't appear.
It is the right error about $\gref$ because $\gref$ is exactly what $\hat\vG$ estimates; so
$\hat\vG - \gref \sim \mathcal N(0, \Ftot^{-1})$, and the pulls about $\gref$ are standard normal.
(c) $w_1 = 2500$, $w_2 = 278$: $g_{\rm ref} = 50/2778 = 0.0180$; plain mean $0.01$. The stack estimates $0.0180$,
with error $0.019$. Comparing it with $0.01$ would call a correct stack biased.
(d) The cell-centre shear is not what the stack converges to. The difference is a deterministic binning effect,
largest where the shear changes quickly. (Full derivation: \texttt{CELL\_REFERENCE\_SHEAR.md}, Secs.~3--5.)
}

% --------------------------------------------------------------------
\problem{B5}{\py\ The tangential frame, and checking a stack with pulls}

\begin{given}
About a lens at the origin, a galaxy at position angle $\phi$ has tangential and cross shear
$g_t = -\mathrm{Re}(g\,e^{-2i\phi})$, $g_\times = -\mathrm{Im}(g\,e^{-2i\phi})$, with $g = g_1 + i g_2$.
\source{\cBStan.}
As a matrix, $(g_t, g_\times)^\T = R(\phi)\,(g_1, g_2)^\T$ with
$R(\phi) = -\begin{pmatrix}\cos2\phi & \sin2\phi\\ -\sin2\phi & \cos2\phi\end{pmatrix}$.
\end{given}

\need{A3, B1, B4.}

\begin{enumerate}
  \item Show that $R(\phi)$ is orthogonal ($RR^\T = I$), and that a likelihood with precision $F$ in the image frame has
    precision $F' = RFR^\T$ in the tangential frame.
  \item An axisymmetric lens produces purely tangential shear: $g_\times = 0$ and $g_t > 0$. Write $(g_1, g_2)$ at angle
    $\phi$ in terms of $g_t$.
  \item \py\ Simulate a cell of 5 galaxies at different $\phi$ around a lens with $g_t$ varying from 0.01 to 0.03 and
    anisotropic, randomly oriented likelihood covariances. Draw noisy measurements, rotate each into the tangential
    frame, stack with A3, and repeat 2000 times. Show that the pulls about $\gref$ have mean $\approx 0$ and spread
    $\approx 1$, while the pulls about the plain mean of the $g_{t,i}$ are biased.
\end{enumerate}

\solution{
(a) Writing $R = -Q$ with $Q$ a rotation by $2\phi$, $RR^\T = QQ^\T = I$. If
$\ln L \propto -\tfrac12 (\vg - \hat\vg)^\T F(\vg - \hat\vg)$ and $\vg = R^\T\vg'$, then
$(\vg - \hat\vg)^\T F(\vg - \hat\vg) = (\vg' - R\hat\vg)^\T RFR^\T(\vg' - R\hat\vg)$, so $F' = RFR^\T$.
(b) $g\,e^{-2i\phi} = -g_t$ gives $g = -g_t e^{2i\phi}$: $g_1 = -g_t\cos2\phi$, $g_2 = -g_t\sin2\phi$.
(c) See the solutions notebook: the pulls about $\gref$ have mean and spread close to $0$ and $1$. The pulls about
the plain mean are offset (by about $0.05$--$0.1\sigma$ in this example; more when the precisions differ more), because
precise galaxies at particular $g_t$ dominate the stack.
}
